\akhead{Topic 5 — Linear Inequalities in One or Two Variables}

\ans{1}{1}{A}{Subtract $9$ from both sides. The direction of the sign is
unchanged by adding or subtracting, so $x \le 6$.}
\ans{1}{2}{8}{Divide both sides by the positive number $5$: $x \le 8$, so the
greatest value is $8$.}
\ans{1}{3}{A}{Divide both sides by $-3$. Dividing by a negative number
reverses the sign, so $x < -4$.}
\ans{1}{4}{C}{Add $3$: $2x > 10$, so $x > 5$. The least integer strictly
greater than $5$ is $6$.}
\ans{1}{5}{B}{If $n$ is the number of notebooks, $6n \le 40$, so
$n \le 6.\overline{6}$. A whole notebook count rounds inward to $6$.}
\ans{1}{6}{16}{Subtract $3$: $\dfrac{x}{4} \ge 4$; multiply by $4$:
$x \ge 16$, so the least value is $16$.}
\ans{1}{7}{C}{Subtract $1$ from all three parts: $-3 < x \le 4$. Of the
options only $4$ lies in that interval; $-3$ is excluded by the strict sign.}
\ans{1}{8}{C}{Subtract $3x$ and add $8$: $10 \le 2x$, so $x \ge 5$ and the
least value is $5$.}
\ans{1}{9}{B}{Distribute: $4x - 12 < 2x + 6$, so $2x < 18$ and $x < 9$. The
greatest integer strictly less than $9$ is $8$.}
\ans{1}{10}{25}{With $n$ crates, $1200 + 45n \le 2325$, so
$45n \le 1125$ and $n \le 25$.}
\ans{1}{11}{B}{Add $1$ to all parts: $-6 \le 2x < 10$; divide by $2$:
$-3 \le x < 5$. The integers $-3$ through $4$ number $8$.}
\ans{1}{12}{A}{Subtract $2$: $-\dfrac{x}{3} \ge 3$; multiply by $-3$ and
reverse the sign: $x \le -9$.}
\ans{1}{13}{60}{Subtract $\dfrac{x}{4}$ and add $3$: $15 \le \dfrac{x}{4}$,
so $x \ge 60$ and the least value is $60$.}
\ans{1}{14}{B}{With $m$ minutes, $25 + 0.10m \le 60$, so $0.10m \le 35$ and
$m \le 350$.}
\ans{1}{15}{B}{At $(0,0)$: $0 > 2(0) - 1 = -1$ is true and
$0 \le -0 + 5 = 5$ is true. Each other pair fails at least one inequality.}
\ans{1}{16}{A}{The term $15h$ grows by $15$ for each additional hour $h$, so
$15$ is the charge per hour; the $45$ is the one-time fee.}
\ans{1}{17}{4}{Add $2x$ and $7$: $12 \ge 3x$, so $x \le 4$ and the greatest
value is $4$.}
\ans{1}{18}{B}{Subtract $k$ and divide by $3$: $x > \dfrac{12 - k}{3}$.
Matching $x > 5$ gives $12 - k = 15$, so $k = -3$.}
\ans{1}{19}{B}{The two boundary lines are parallel, so the system is empty
exactly when $4x + c$ stays below $4x + 7$, that is $c < 7$. The greatest
integer less than $7$ is $6$.}
\ans{1}{20}{6}{Subtract $3$: $-7 \le -2x < 6$; divide by $-2$ and reverse:
$3.5 \ge x > -3$. The integers $-2$ through $3$ number $6$.}
\ans{1}{21}{B}{With $x = 4$: $12 + 2y \le 30$ gives $y \le 9$, and
$4 + y \ge 12$ gives $y \ge 8$. The greatest value allowed is $9$.}
\ans{1}{22}{5}{Subtract $3x$ and $8$: $(a - 3)x \le 12$. For the solution to
read $x \le 6$ the coefficient $a - 3$ must be positive with
$\dfrac{12}{a-3} = 6$, so $a - 3 = 2$ and $a = 5$.}

\ans{2}{1}{A}{Distribute: $6 - 2x > 4x - 6$, so $12 > 6x$ and $x < 2$. The
greatest integer strictly less than $2$ is $1$.}
\ans{2}{2}{A}{Subtract $4$: $-\dfrac{2}{3}x \le 6$; multiply by
$-\dfrac{3}{2}$ and reverse the sign: $x \ge -9$.}
\ans{2}{3}{50}{Solve $\frac{9}{5}t + 32 \le 122$: $\frac{9}{5}t \le 90$, so
$t \le 50$. (The lower bound gives $t \ge 20$.)}
\ans{2}{4}{B}{At $(1,2)$: $2 < -2(1) + 6 = 4$ is true and
$2 \ge 1 - 3 = -2$ is true. Every other pair fails one of the two.}
\ans{2}{5}{B}{With $m$ miles, $2.50 + 1.75m \le 30$, so $1.75m \le 27.50$ and
$m \le 15.71\ldots$; a whole number of miles rounds inward to $15$.}
\ans{2}{6}{5}{Multiply all parts by $2$: $-2 \le 3x - 4 < 14$; add $4$:
$2 \le 3x < 18$; divide by $3$: $\frac{2}{3} \le x < 6$. The greatest integer
below $6$ is $5$.}
\ans{2}{7}{C}{Add $9$: $kx < 15$. Dividing by a positive $k$ keeps the sign,
so $x < \dfrac{15}{k}$. Matching $x < 5$ gives $k = 3$.}
\ans{2}{8}{A}{At most $15$ trays gives $s + \ell \le 15$; the servings
$8s + 20\ell$ must be at least $200$, giving $8s + 20\ell \ge 200$.}
\ans{2}{9}{6}{With $x = 3$: $36 + 7y \le 84$, so $7y \le 48$ and
$y \le 6.85\ldots$; the greatest integer is $6$.}
\ans{2}{10}{A}{Each extra day subtracts another $18$ from the $250$ liters,
so $18$ is the number of liters used per day.}
\ans{2}{11}{C}{Multiply by $6$: $2(x+2) - 3(x-4) \ge 6$, so
$-x + 16 \ge 6$. Then $-x \ge -10$, and reversing gives $x \le 10$.}
\ans{2}{12}{10}{The perimeter is $2\bigl(w + (w+3)\bigr) = 4w + 6 \le 46$, so
$4w \le 40$ and $w \le 10$.}
\ans{2}{13}{C}{Distribute: $-3x + 12 \ge 2x + 7$, so $5 \ge 5x$ and
$x \le 1$; the greatest value is $1$.}
\ans{2}{14}{A}{Substitute $y = x - 3$ into $3x + y = 21$: $4x - 3 = 21$, so
$x = 6$ and $y = 6 - 3 = 3$.}
\ans{2}{15}{B}{Subtract $5$: $-\dfrac{x-1}{4} > -3$; multiply by $-4$ and
reverse: $x - 1 < 12$, so $x < 13$ and the greatest integer is $12$.}
\ans{2}{16}{5}{Subtract $2x$ and add $4$: $c + 4 \le 3x$, so
$x \ge \dfrac{c+4}{3}$. Matching $x \ge 3$ gives $c + 4 = 9$, so $c = 5$.}
\ans{2}{17}{B}{From $4 - 3x > -2$: $x < 2$. From $4 - 3x \le 13$: $x \ge -3$.
So $-3 \le x < 2$, which contains the $5$ integers $-3$ through $1$.}
\ans{2}{18}{C}{From $2 < x < 6$: $4 < 2x < 12$, so $-12 < -2x < -4$ and
$-7 < 5 - 2x < 1$. The greatest integer in that range is $0$.}
\ans{2}{19}{12}{Deluxe chairs are maximized when labor spent on standard
chairs is least, at $10$ standard chairs using $20$ hours. Then
$5y \le 80 - 20 = 60$, so $y \le 12$.}
\ans{2}{20}{A}{Multiplying both sides of $a < b$ by the negative number $c$
reverses the sign, giving $ac > bc$.}
\ans{2}{21}{A}{Subtract $6$: $mx > 12$. Dividing by the negative constant $m$
reverses the sign, giving $x < \dfrac{12}{m}$.}
\ans{2}{22}{8}{Both upper bounds cap $y$, and the cap is largest where the
lines meet: $-2x + 14 = 3x - 1$ gives $x = 3$, so $y = 3(3) - 1 = 8$.}
